% ===========================================================================
%  Chapitre 13 — Probabilités
%  PE 2026, composante TRAITEMENT DES DONNÉES — série OSE
% ===========================================================================
\bmtchapitre{Probabilités}
  {Résoudre des problèmes relatifs aux probabilités : vocabulaire des
   événements, calcul dans le cas d'équiprobabilité, tirages successifs et
   simultanés, probabilité conditionnelle et indépendance.}
  {Traitement des données \textbullet\ environ 12 heures}

Le hasard se calcule. C'est une affirmation surprenante, et pourtant tout ce
chapitre repose sur elle : dès qu'on sait décrire l'ensemble de tout ce qui
peut arriver, on sait mesurer la part qu'occupe chaque événement. La
probabilité n'est rien d'autre que cette mesure, comprise entre zéro et un —
et c'est exactement l'outil qu'utilise un assureur pour fixer une prime, ou
une coopérative pour estimer un risque de récolte.

Le chapitre se joue en deux temps. Le premier est un travail de description :
énumérer, dénombrer, distinguer un tirage avec remise d'un tirage sans remise.
Le second introduit le conditionnement — que devient une probabilité quand on
sait déjà quelque chose ? — et c'est là que se trouvent les questions les plus
subtiles, celles où l'intuition trompe régulièrement.

% ---------------------------------------------------------------------------
\begin{revision}
Dénombrement et probabilités de première.

\begin{enumerate}[label=\textbf{\arabic*.}, leftmargin=*, itemsep=0.35em]
  \item Combien y a-t-il de façons de choisir $2$ objets parmi $5$, sans tenir
        compte de l'ordre ?
  \item Calculer $\combi{2}{5}$, $\combi{3}{6}$ et $\arrang{2}{5}$.
  \item Un dé équilibré est lancé une fois : quelle est la probabilité
        d'obtenir un nombre pair ?
  \item Que vaut $\proba{\evtc{A}}$ en fonction de $\proba{A}$ ?
  \item Une urne contient $4$ boules rouges et $6$ vertes. On en tire une :
        probabilité qu'elle soit rouge ?
\end{enumerate}
\end{revision}

\begin{automatismes}
Sans calculatrice, sans brouillon.

\begin{enumerate}[label=\textbf{\alph*)}, leftmargin=*, itemsep=0.25em]
  \item $\combi{2}{4}$
  \item $\combi{1}{7}$
  \item $\combi{0}{5}$
  \item $\arrang{2}{4}$
  \item $\proba{A} = 0{,}3$ : que vaut $\proba{\evtc{A}}$ ?
  \item Probabilité d'obtenir « pile » deux fois de suite
  \item Nombre de tirages successifs avec remise de $2$ boules parmi $5$
  \item Nombre de tirages simultanés de $2$ boules parmi $5$
\end{enumerate}
\end{automatismes}

\begin{corrige}[automatismes]
a) $6$ \quad b) $7$ \quad c) $1$ \quad d) $12$ \quad e) $0{,}7$ \quad
f) $\frac14$ \quad g) $25$ \quad h) $10$.
\end{corrige}

% ===========================================================================
\section{Vocabulaire}

\begin{definition}[univers, événement]
L'\textbf{univers} $\Omega$ d'une expérience aléatoire est l'ensemble de tous
ses résultats possibles ; chacun d'eux est une \textbf{éventualité}. Un
\textbf{événement} est une partie de $\Omega$. L'événement $\Omega$ est
\textbf{certain}, l'ensemble vide est \textbf{impossible}.

L'événement \textbf{contraire} de $A$, noté $\evtc{A}$, est l'ensemble des
éventualités qui ne sont pas dans $A$. Deux événements sont
\textbf{incompatibles} lorsque leur intersection est vide.
\end{definition}

\begin{avertissement}
« Contraire » et « incompatible » ne sont pas synonymes. Deux événements
contraires sont incompatibles et leur réunion est $\Omega$ ; deux événements
incompatibles peuvent laisser de la place à un troisième. Sur un dé, « obtenir
$1$ » et « obtenir $2$ » sont incompatibles sans être contraires.
\end{avertissement}

% ===========================================================================
\section{Probabilité et équiprobabilité}

\begin{definition}[loi de probabilité]
Une \textbf{probabilité} sur un univers fini $\Omega$ est une application
$\Prob$ qui à chaque événement associe un réel de $\intervalle{0}{1}$, telle
que $\proba{\Omega} = 1$ et que $\proba{A \cup B} = \proba{A}+\proba{B}$ dès
que $A$ et $B$ sont incompatibles.
\end{definition}

\begin{propriete}[cas d'équiprobabilité]
Lorsque toutes les éventualités ont la même probabilité,
\[
  \proba{A} = \frac{\text{nombre de cas favorables}}{
  \text{nombre de cas possibles}} = \frac{\operatorname{card}A}{
  \operatorname{card}\Omega}.
\]
\end{propriete}

\begin{avertissement}
Cette formule n'est valable que dans le cas d'équiprobabilité, et l'énoncé
doit le justifier : dé équilibré, boules indiscernables au toucher, tirage au
hasard. Un dé pipé ou une urne dont les boules ne sont pas indiscernables
interdit son emploi.
\end{avertissement}

\begin{exemple}[un dé déséquilibré]
Un dé cubique porte les numéros $0$, $1$, $1$, $2$, $2$, $2$. Les faces sont
équiprobables, mais pas les numéros :
\[
  \proba{0} = \frac16, \qquad \proba{1} = \frac26 = \frac13,
  \qquad \proba{2} = \frac36 = \frac12 .
\]
C'est bien l'équiprobabilité des \emph{faces} qui permet ce calcul ; les
numéros, eux, héritent de probabilités différentes.
\end{exemple}

% ===========================================================================
\section{Réunion, intersection, contraire}

\begin{propriete}[formules de base]
Pour tous événements $A$ et $B$,
\[
  \proba{\evtc{A}} = 1-\proba{A},
  \qquad
  \proba{A \cup B} = \proba{A}+\proba{B}-\proba{A \cap B}.
\]
\end{propriete}

\begin{methode}{« au moins un » se traite par le contraire}
Quand un énoncé demande la probabilité d'obtenir \emph{au moins une fois} un
résultat, on calcule d'abord la probabilité de ne l'obtenir \emph{jamais} —
un seul produit — puis on soustrait à $1$. C'est la question la plus fréquente
de l'exercice de probabilités au baccalauréat, et le calcul direct y est
toujours plus long.
\end{methode}

% ===========================================================================
\section{Tirages successifs et tirages simultanés}

\begin{propriete}[trois protocoles à ne pas confondre]
\begin{center}
\setlength{\tabcolsep}{7pt}
\renewcommand{\arraystretch}{1.3}
\begin{tabular}{@{}lll@{}}
\toprule
\textsf{\footnotesize\bfseries PROTOCOLE} &
\textsf{\footnotesize\bfseries L'ORDRE COMPTE} &
\textsf{\footnotesize\bfseries NOMBRE DE TIRAGES} \\
\midrule
Successif avec remise & oui & $n^{p}$ \\
Successif sans remise & oui & $\arrang{p}{n}$ \\
Simultané & non & $\combi{p}{n}$ \\
\bottomrule
\end{tabular}
\end{center}
\end{propriete}

\begin{remarque}
Un tirage simultané de $p$ boules donne les mêmes probabilités qu'un tirage
successif sans remise dont on oublie l'ordre. On peut donc traiter un tirage
simultané par un arbre, à condition de sommer toutes les branches conduisant
au même contenu final.
\end{remarque}

\begin{exemple}[avec ou sans remise]
Une urne contient trois boules blanches et sept noires. On tire deux boules.

\emph{Sans remise.} $\proba{\text{deux noires}}
= \frac{7}{10} \times \frac69 = \frac{42}{90} = \frac{7}{15}$.

\emph{Avec remise.} $\proba{\text{deux noires}}
= \frac{7}{10} \times \frac{7}{10} = \frac{49}{100}$.

Les deux résultats diffèrent : le premier tirage modifie la composition de
l'urne dans le premier cas, pas dans le second. Lire attentivement le
protocole est donc la première chose à faire.
\end{exemple}

\begin{entrainement}[dénombrer et calculer]
\exo[\niveau{1}] Une urne contient $4$ boules rouges et $6$ vertes. On tire
deux boules simultanément. Probabilité d'obtenir deux rouges ?

\exo[\niveau{1}] Même urne, tirages successifs avec remise. Même question.

\exo[\niveau{2}] On lance trois fois un dé équilibré. Probabilité d'obtenir au
moins un $6$ ?

\exo[\niveau{2}] Un dé porte les numéros $2$, $3$, $3$, $3$, $4$, $4$. On le
lance trois fois. Probabilité que les trois numéros soient premiers ?

\exo[\niveau{3}] Une urne contient $9$ boules : quatre portent le numéro $1$,
trois le numéro $2$, deux le numéro $3$. On tire deux boules simultanément et
l'on note $b$ la somme. Donner la loi de $b$.
\end{entrainement}

% ===========================================================================
\section{Arbres de probabilité}

\begin{propriete}[règles de l'arbre]
\begin{enumerate}[label=\textbf{\arabic*.}, leftmargin=*, itemsep=0.15em]
  \item La somme des probabilités portées par les branches issues d'un même
        n\oe ud vaut $1$.
  \item La probabilité d'un chemin est le produit des probabilités de ses
        branches.
  \item La probabilité d'un événement est la somme des probabilités des
        chemins qui le réalisent.
\end{enumerate}
\end{propriete}

\begin{visualisation}[un arbre à deux niveaux]
\begin{center}
\begin{tikzpicture}[x=1cm, y=1cm, font=\footnotesize,
  level 1/.style={sibling distance=2.6cm},
  level 2/.style={sibling distance=1.35cm}]
  \node[bmtEncre] {}
    child { node[bmtIndigo] {$A$}
      child { node[bmtGris] {$B$} edge from parent
        node[left, bmtGris, pos=0.55] {$\probacond{B}{A}$} }
      child { node[bmtGris] {$\evtc{B}$} }
      edge from parent node[left, bmtIndigo, pos=0.55] {$\proba{A}$} }
    child { node[bmtLaterite] {$\evtc{A}$}
      child { node[bmtGris] {$B$} }
      child { node[bmtGris] {$\evtc{B}$} }
      edge from parent node[right, bmtLaterite, pos=0.55]
        {$\proba{\evtc{A}}$} };
\end{tikzpicture}
\end{center}

Chaque chemin complet représente une histoire possible, et sa probabilité est
le produit des nombres rencontrés en descendant. Le premier étage donne les
probabilités simples, le second des probabilités \emph{conditionnelles} : la
branche $\probacond{B}{A}$ se lit « probabilité de $B$ sachant que $A$ est
déjà réalisé ». C'est l'arbre qui rend cette notion naturelle, bien avant sa
définition formelle.
\end{visualisation}

% ===========================================================================
\section{Probabilité conditionnelle}

\begin{definition}[probabilité conditionnelle]
Soient $A$ et $B$ deux événements avec $\proba{A} \neq 0$. La
\textbf{probabilité de $B$ sachant $A$} est
\[
  \probacond{B}{A} = \frac{\proba{A \cap B}}{\proba{A}} .
\]
\end{definition}

\begin{propriete}[formule des probabilités totales]
Si $A_{1},\dots,A_{n}$ forment une partition de $\Omega$ en événements de
probabilité non nulle, alors pour tout événement $B$,
\[
  \proba{B} = \sum_{i=1}^{n}\proba{A_{i}} \times \probacond{B}{A_{i}} .
\]
En particulier
$\proba{B} = \proba{A}\probacond{B}{A}
+\proba{\evtc{A}}\probacond{B}{\evtc{A}}$.
\end{propriete}

\begin{avertissement}
$\probacond{B}{A}$ et $\probacond{A}{B}$ n'ont aucune raison d'être égales.
La probabilité qu'un producteur soit contrôlé positif sachant qu'il a fraudé,
et celle qu'il ait fraudé sachant qu'il est contrôlé positif, sont deux
nombres différents. Confondre les deux est l'erreur la plus coûteuse du
chapitre.
\end{avertissement}

\begin{exemple}[remonter l'arbre]
Une enquête indique que $75\,\%$ des personnes interrogées croient au
réchauffement. Parmi elles, $60\,\%$ se disent écologistes ; parmi celles qui
n'y croient pas, $8\,\%$ ne se disent pas écologistes.

On note $R$ « croire au réchauffement » et $E$ « se dire écologiste ». Alors
$\proba{R} = 0{,}75$, $\probacond{E}{R} = 0{,}6$ et
$\probacond{\evtc{E}}{\evtc{R}} = 0{,}08$, donc
\[
  \proba{\evtc{R} \cap E} = 0{,}25 \times 0{,}92 = 0{,}23,
  \qquad
  \proba{E} = 0{,}75 \times 0{,}6+0{,}23 = 0{,}68 .
\]
La probabilité qu'une personne écologiste ne croie pas au réchauffement vaut
alors $\dfrac{0{,}23}{0{,}68} \approx 0{,}34$ : un tiers, ce qui est loin
d'être négligeable et n'apparaissait nulle part dans les données de départ.
\end{exemple}

\begin{entrainement}[conditionnement]
\exo[\niveau{1}] $\proba{A} = 0{,}4$ et $\proba{A \cap B} = 0{,}1$.
Calculer $\probacond{B}{A}$.

\exo[\niveau{2}] Deux urnes : la première contient $3$ boules blanches et $1$
noire, la seconde $1$ blanche et $1$ noire. On tire une boule de la première
et on la met dans la seconde, puis on tire une boule de la seconde.
Probabilité qu'elle soit blanche ?

\exo[\niveau{2}] Dans une population, $2\,\%$ des personnes sont porteuses
d'un caractère. Un test détecte le caractère dans $95\,\%$ des cas s'il est
présent, et le signale à tort dans $3\,\%$ des cas sinon. Une personne est
testée positive : quelle est la probabilité qu'elle soit porteuse ?

\exo[\niveau{3}] Deux chasseurs tirent simultanément. Le premier atteint sa
cible trois fois sur quatre, le second trois fois sur cinq. Probabilité que le
gibier soit touché ?
\end{entrainement}

% ===========================================================================
\section{Événements indépendants}

\begin{definition}[indépendance]
Deux événements $A$ et $B$ sont \textbf{indépendants} lorsque
\[
  \proba{A \cap B} = \proba{A} \times \proba{B},
\]
ce qui équivaut, si $\proba{A} \neq 0$, à
$\probacond{B}{A} = \proba{B}$ : savoir que $A$ est réalisé ne change rien à
la probabilité de $B$.
\end{definition}

\begin{avertissement}
Indépendants et incompatibles sont des notions opposées, non voisines. Deux
événements incompatibles de probabilités non nulles ne sont \emph{jamais}
indépendants, puisque $\proba{A \cap B} = 0$ tandis que le produit des
probabilités ne l'est pas.
\end{avertissement}

\begin{propriete}[répétition d'épreuves indépendantes]
Si l'on répète $n$ fois, de façon indépendante, une épreuve où un événement
$A$ a la probabilité $p$, alors
\[
  \proba{\text{$A$ n'est jamais réalisé}} = (1-p)^{n},
\]
\[
  \proba{\text{$A$ est réalisé au moins une fois}} = 1-(1-p)^{n}.
\]
\end{propriete}

\begin{exemple}[combien de répétitions ?]
Un événement de probabilité $\frac14$ est répété $n$ fois de façon
indépendante. On veut $p_{n} = 1-\left(\frac34\right)^{n} \geq 0{,}99$.

Cela équivaut à $\left(\frac34\right)^{n} \leq 0{,}01$, soit
\[
  n \geq \frac{\ln 0{,}01}{\ln 0{,}75} \approx 16{,}01 .
\]
Le plus petit entier convenable est donc $n = 17$.
\end{exemple}

\begin{entrainement}[indépendance]
\exo[\niveau{1}] $\proba{A} = 0{,}5$, $\proba{B} = 0{,}4$ et
$\proba{A \cap B} = 0{,}2$. Les événements sont-ils indépendants ?

\exo[\niveau{2}] On lance $n$ fois un dé équilibré. À partir de quel $n$ la
probabilité d'obtenir au moins un $6$ dépasse-t-elle $0{,}9$ ?

\exo[\niveau{3}] On répète $n$ fois une épreuve où un événement a la
probabilité $\frac13$. Déterminer $n$ pour que la probabilité de ne jamais le
réaliser soit inférieure à $10^{-3}$.
\end{entrainement}

% ===========================================================================
\begin{annales}
Les probabilités occupent le second volet de l'exercice 1 de l'épreuve, et le
scénario est stable : une urne ou un dé, quelques probabilités simples, puis
une répétition d'épreuves conduisant à une inéquation. Les énoncés qui suivent
sont repris de sujets réellement posés.

\exoann{Baccalauréat, Madagascar, série S, session 2021 — exercice 1, partie II}
On dispose de quatre trous $T_{1}$, $T_{2}$, $T_{3}$, $T_{4}$ et de quatre
billes numérotées $1$, $2$, $3$ et $4$. Chaque bille entre dans un trou et
chaque trou peut recevoir zéro à quatre billes ; tous les événements
élémentaires sont équiprobables. On lance une fois les quatre billes. Soit $A$
l'événement « la bille numéro $2$ est dans le trou $T_{2}$ ». Montrer que
$\proba{A} = \frac14$.

\exoann{Baccalauréat, Madagascar, série S, session 2021 — exercice 1, partie II, suite}
On lance $n$ fois de suite, de manière indépendante, les quatre billes vers
les quatre trous. On note $p_{n}$ la probabilité de l'événement $A_{n}$ :
« $A$ est réalisé au moins une fois lors de ces $n$ lancers ». Calculer
$p_{n}$ en fonction de $n$, puis déterminer le plus petit entier naturel $n$
tel que $p_{n} \geq 0{,}99$.

\exoann{Baccalauréat, Madagascar, série S, session 2023 — exercice 1, partie B}
On dispose de deux urnes $U_{1}$ et $U_{2}$. L'urne $U_{1}$ contient trois
boules blanches et une boule noire, l'urne $U_{2}$ une blanche et une noire ;
toutes les boules sont indiscernables au toucher. L'épreuve $E$ consiste à
tirer une boule de $U_{1}$, à la mettre dans $U_{2}$, puis à tirer au hasard et
simultanément deux boules de $U_{2}$. Soit $A$ l'événement « obtenir deux
boules blanches au deuxième tirage ». Montrer que
$\proba{A} = \frac14$.

\exoann{Baccalauréat blanc, lycée Philibert Tsiranana, Terminale S, 2022-2023 — exercice 1, partie II}
On dispose d'un dé cubique équilibré dont les faces sont numérotées de $1$ à
$6$, et d'une urne contenant quatre boules numérotées $1$, trois numérotées
$2$ et deux numérotées $3$. On lance le dé et l'on note $a$ le numéro obtenu ;
on tire simultanément deux boules de l'urne et l'on note $b$ la somme de leurs
numéros. Calculer la probabilité des événements
\[
  A : \text{« } ab = 8 \text{ »}, \qquad
  B : \text{« } a+b<5 \text{ »} .
\]

\exoann{Baccalauréat blanc, lycée Philibert Tsiranana, Terminale S, 2022-2023 — exercice 1, partie II, suite}
On lance $n$ fois de suite et de manière indépendante ce même dé. Soit
$A_{n}$ l'événement « ne jamais obtenir un numéro multiple de trois au cours
de ces $n$ lancers » et $p_{n} = \proba{A_{n}}$. Déterminer $p_{n}$, puis
calculer $\lim_{n\to+\infty}p_{n}$.

\exoann{Baccalauréat, Congo-Brazzaville, série C, session 2023 — exercice 4}
Une urne contient dix boules : trois blanches et sept noires. On effectue deux
tirages successifs sans remise, et l'on note $X$ le nombre de boules blanches
obtenues.
\begin{enumerate}[label=\textbf{\alph*)}, leftmargin=*, itemsep=0.15em]
  \item Construire l'arbre de probabilités.
  \item Déterminer $X(\Omega)$ et démontrer que
        $\proba{X = 0} = \frac{7}{15}$.
  \item Donner la loi de probabilité de $X$.
\end{enumerate}

\exoann{Baccalauréat, Côte d'Ivoire, série C, session 2023 — exercice 3}
Une enquête révèle que $75\,\%$ des scientifiques interrogés croient au
réchauffement climatique. Parmi eux, $60\,\%$ se déclarent écologistes ; parmi
ceux qui n'y croient pas, $8\,\%$ ne se déclarent pas écologistes. On note $R$
et $E$ les événements correspondants.
\begin{enumerate}[label=\textbf{\alph*)}, leftmargin=*, itemsep=0.15em]
  \item Construire un arbre pondéré, puis donner $\proba{R}$,
        $\probacond{\evtc{E}}{\evtc{R}}$ et $\probacond{E}{R}$.
  \item Justifier que $\proba{\evtc{R} \cap E} = 0{,}23$, puis que
        $\proba{E} = 0{,}68$.
  \item Calculer la probabilité qu'un écologiste ne croie pas au
        réchauffement, arrondie au centième.
\end{enumerate}

\exoann{Baccalauréat, Madagascar, série C, session 2011 — exercice, partie A}
On dispose d'un dé cubique parfait dont les faces sont numérotées $0$, $1$,
$1$, $2$, $2$, $2$.
\begin{enumerate}[label=\textbf{\alph*)}, leftmargin=*, itemsep=0.15em]
  \item Calculer les probabilités $P_{0}$, $P_{1}$, $P_{2}$ d'apparition
        respective des numéros $0$, $1$ et $2$ en un lancer.
  \item On lance trois fois ce dé de manière indépendante. Calculer la
        probabilité de « la somme des numéros obtenus vaut $4$ », puis de
        « obtenir exactement deux fois le numéro $1$ », puis de « obtenir
        trois numéros différents ».
\end{enumerate}

\exoann[arbre pondéré et probabilité conditionnelle]{Baccalauréat, France, série ES/L, Liban, 27 mai 2014 — exercice 1, partie A (extrait)}
Un serveur, travaillant dans une pizzeria, remarque qu'en moyenne $40\,\%$
des clients sont des familles, $25\,\%$ des personnes seules et $35\,\%$ des
couples. Il note aussi que $70\,\%$ des familles laissent un pourboire,
$90\,\%$ des personnes seules en laissent un, et $40\,\%$ des couples en
laissent un. On prend au hasard une table occupée et l'on considère les
événements $F$ « la table est occupée par une famille », $S$ « \ldots\ par
une personne seule », $C$ « \ldots\ par un couple » et $R$ « le serveur
reçoit un pourboire ».
\begin{enumerate}[label=\textbf{\arabic*.}, leftmargin=*, itemsep=0.2em]
  \item Préciser les probabilités $\proba{F}$ et $\probacond{R}{S}$.
  \item Construire l'arbre pondéré correspondant à la situation.
  \item Calculer $\proba{F\cap R}$, puis $\proba{R}$.
  \item Sachant que le serveur a reçu un pourboire, calculer la probabilité
        que ce pourboire vienne d'un couple (résultat arrondi à $10^{-3}$).
\end{enumerate}
\end{annales}

\begin{corrige}[annales]
\textbf{1.} La bille numéro $2$ choisit son trou parmi quatre, tous
équiprobables, et son choix ne dépend pas de celui des autres billes :
$\proba{A} = \frac14$.

\textbf{2.} Les lancers étant indépendants,
$\proba{\evtc{A_{n}}} = \left(\frac34\right)^{n}$, donc
$p_{n} = 1-\left(\frac34\right)^{n}$. L'inéquation
$\left(\frac34\right)^{n} \leq 0{,}01$ donne
$n \geq \frac{\ln0{,}01}{\ln0{,}75} \approx 16{,}01$ : le plus petit entier
est $n = 17$.

\textbf{3.} Si la boule transférée est blanche — probabilité $\frac34$ —,
$U_{2}$ contient deux blanches et une noire, et la probabilité de tirer les
deux blanches vaut $\frac{\combi{2}{2}}{\combi{2}{3}} = \frac13$. Si elle est
noire — probabilité $\frac14$ —, $U_{2}$ ne contient qu'une blanche et cette
probabilité est nulle. Donc
$\proba{A} = \frac34 \times \frac13 = \frac14$.

\textbf{4.} Les sommes possibles et leurs probabilités, obtenues par tirage
simultané de deux boules parmi neuf, sont
$\proba{b = 2} = \frac{6}{36} = \frac16$,
$\proba{b = 3} = \frac{12}{36} = \frac13$,
$\proba{b = 4} = \frac{11}{36}$,
$\proba{b = 5} = \frac{6}{36} = \frac16$,
$\proba{b = 6} = \frac{1}{36}$. Alors
$\proba{A} = \frac16 \times \frac{11}{36}+\frac16 \times \frac16
= \frac{17}{216}$, et
$\proba{B} = \frac16\left(\frac16+\frac16+\frac13\right) = \frac19$.

\textbf{5.} Les multiples de trois sur le dé sont $3$ et $6$, de probabilité
$\frac13$. Donc $p_{n} = \left(\frac23\right)^{n}$, qui tend vers $0$.

\textbf{6.} L'arbre a deux niveaux, les probabilités du second dépendant du
résultat du premier. $X(\Omega) = \{0\,;1\,;2\}$ et
$\proba{X = 0} = \frac{7}{10} \times \frac69 = \frac{7}{15}$. De même
$\proba{X = 1} = 2 \times \frac{3}{10} \times \frac79 = \frac{7}{15}$ et
$\proba{X = 2} = \frac{3}{10} \times \frac29 = \frac{1}{15}$ ; la somme vaut
bien $1$.

\textbf{7.} $\proba{R} = 0{,}75$, $\probacond{\evtc{E}}{\evtc{R}} = 0{,}08$,
$\probacond{E}{R} = 0{,}6$. Alors
$\proba{\evtc{R} \cap E} = 0{,}25 \times 0{,}92 = 0{,}23$ et
$\proba{E} = 0{,}75 \times 0{,}6+0{,}23 = 0{,}68$. Enfin
$\probacond{\evtc{R}}{E} = \frac{0{,}23}{0{,}68} \approx 0{,}34$.

\textbf{8. a)} $P_{0} = \frac16$, $P_{1} = \frac13$, $P_{2} = \frac12$.
\textbf{b)} Somme égale à $4$ : les répartitions possibles sont
$\{0\,;2\,;2\}$ et $\{1\,;1\,;2\}$, de probabilités respectives
$3 \times \frac16 \times \frac12 \times \frac12 = \frac18$ et
$3 \times \frac13 \times \frac13 \times \frac12 = \frac16$ ; le total vaut
$\frac{7}{24}$. Exactement deux fois le $1$ :
$\combi{2}{3}\left(\frac13\right)^{2}\frac23 = \frac29$. Trois numéros
différents : $3! \times \frac16 \times \frac13 \times \frac12 = \frac16$.

\textbf{9. 1.} $\proba{F} = 0{,}40$ (donnée de l'énoncé) et
$\probacond{R}{S} = 0{,}90$ (les personnes seules laissent un pourboire neuf
fois sur dix).
\textbf{2.} L'arbre part de $F$ ($0{,}40$), $S$ ($0{,}25$) et $C$ ($0{,}35$) ;
chacune de ces trois branches se prolonge par $R$ et $\overline{R}$, de
probabilités conditionnelles respectives $0{,}70/0{,}30$ (sachant $F$),
$0{,}90/0{,}10$ (sachant $S$) et $0{,}40/0{,}60$ (sachant $C$).
\textbf{3.} $\proba{F\cap R} = \proba{F}\times\probacond{R}{F}
= 0{,}40\times0{,}70 = 0{,}28$. En sommant sur les trois branches,
\[
  \proba{R} = 0{,}28 + \underbrace{0{,}25\times0{,}90}_{0{,}225}
  + \underbrace{0{,}35\times0{,}40}_{0{,}14} = 0{,}645 .
\]
\textbf{4.} $\proba{C\cap R} = 0{,}35\times0{,}40 = 0{,}14$, donc
\[
  \probacond{C}{R} = \frac{\proba{C\cap R}}{\proba{R}}
  = \frac{0{,}14}{0{,}645} \approx 0{,}217 .
\]
\end{corrige}

% ===========================================================================
\begin{jecherche}[le contrôle des sacs de vanille]
Une coopérative expédie des sacs de vanille. Deux ateliers les préparent :
l'atelier $A$ prépare $60\,\%$ des sacs et l'atelier $B$ le reste. Un sac est
déclaré conforme s'il respecte le calibre. On sait que $95\,\%$ des sacs de
$A$ sont conformes, contre $88\,\%$ de ceux de $B$.

\begin{enumerate}[label=\textbf{\arabic*.}, leftmargin=*, itemsep=0.3em]
  \item Construire un arbre pondéré décrivant la situation.
  \item Calculer la probabilité qu'un sac pris au hasard soit conforme.
  \item Un sac est déclaré non conforme. Quelle est la probabilité qu'il
        vienne de l'atelier $B$ ?
  \item L'inspecteur prélève cinq sacs au hasard, la production étant assez
        grande pour assimiler les prélèvements à des tirages indépendants.
        Quelle est la probabilité qu'ils soient tous conformes ?
  \item Quelle est la probabilité qu'au moins un sac soit non conforme ?
  \item Combien de sacs faut-il prélever pour que la probabilité d'en trouver
        au moins un non conforme dépasse $0{,}95$ ?
\end{enumerate}
\end{jecherche}

\begin{corrige}[le contrôle des sacs de vanille]
\textbf{1.} Premier niveau : $A$ avec $0{,}6$ et $B$ avec $0{,}4$. Second
niveau : conforme avec $0{,}95$ sous $A$ et $0{,}88$ sous $B$.

\textbf{2.} $\proba{C} = 0{,}6 \times 0{,}95+0{,}4 \times 0{,}88
= 0{,}57+0{,}352 = 0{,}922$.

\textbf{3.} $\proba{\evtc{C}} = 0{,}078$ et
$\proba{B \cap \evtc{C}} = 0{,}4 \times 0{,}12 = 0{,}048$, donc
$\probacond{B}{\evtc{C}} = \frac{0{,}048}{0{,}078} \approx 0{,}62$. Un sac
non conforme a donc près de deux chances sur trois de venir de $B$, alors que
$B$ ne produit que $40\,\%$ des sacs.

\textbf{4.} $0{,}922^{5} \approx 0{,}667$.

\textbf{5.} $1-0{,}922^{5} \approx 0{,}333$.

\textbf{6.} Il faut $1-0{,}922^{n} \geq 0{,}95$, soit
$0{,}922^{n} \leq 0{,}05$, c'est-à-dire
$n \geq \frac{\ln0{,}05}{\ln0{,}922} \approx 36{,}9$ : il faut prélever au
moins $37$ sacs.
\end{corrige}

% ===========================================================================
\begin{jemeteste}
Répondre par vrai ou faux, en justifiant en une ligne.

\begin{enumerate}[label=\textbf{\arabic*.}, leftmargin=*, itemsep=0.28em]
  \item Deux événements incompatibles sont indépendants.
  \item $\proba{A \cup B} = \proba{A}+\proba{B}$ pour tous $A$ et $B$.
  \item $\probacond{B}{A} = \probacond{A}{B}$.
  \item Dans un tirage simultané, l'ordre n'intervient pas.
  \item La probabilité d'obtenir au moins un succès en $n$ épreuves vaut
        $1-(1-p)^{n}$.
  \item Un dé pipé permet encore d'utiliser la formule « cas favorables sur
        cas possibles ».
\end{enumerate}
\end{jemeteste}

\begin{corrige}[je me teste]
\textbf{1.} Faux — c'est même le contraire dès que les deux probabilités sont
non nulles.
\textbf{2.} Faux — il faut retrancher $\proba{A \cap B}$, sauf si les
événements sont incompatibles.
\textbf{3.} Faux — les deux quotients ont des dénominateurs différents.
\textbf{4.} Vrai.
\textbf{5.} Vrai, pour des épreuves indépendantes.
\textbf{6.} Faux — cette formule suppose l'équiprobabilité.
\end{corrige}

% ===========================================================================
\begin{commeaubac}
\bmtsource{Baccalauréat, série S, session 2026 — exercice 1, partie II}
Une urne contient dix boules indiscernables au toucher, portant les lettres
$A$, $A$, $A$, $B$, $B$, $B$, $B$, $B$, $C$, $C$. Une épreuve $(E)$ consiste à
tirer au hasard une boule de l'urne.

\begin{enumerate}[label=\textbf{\arabic*.}, leftmargin=*, itemsep=0.25em]
  \item Calculer $\proba{A}$, $\proba{B}$ et $\proba{C}$.
        \hfill\textup{(0,5 pt)}
  \item On répète cinq fois de suite et de manière indépendante l'épreuve
        $(E)$.
        \begin{enumerate}[label=\textbf{\alph*)}, leftmargin=1.4em]
          \item Calculer la probabilité de l'événement $F$ : « la lettre $B$
                apparaît exactement deux fois ». \hfill\textup{(0,5 pt)}
          \item Soit $X$ la variable aléatoire donnant le nombre
                d'apparitions de la lettre $B$ à l'issue des cinq épreuves.
                Calculer $\proba{X \geq 1}$. \hfill\textup{(0,5 pt)}
        \end{enumerate}
\end{enumerate}
\end{commeaubac}

\begin{corrige}[comme au baccalauréat]
\textbf{1.} $\proba{A} = \frac{3}{10}$, $\proba{B} = \frac12$ et
$\proba{C} = \frac15$ ; la somme vaut bien $1$.

\textbf{2. a)} Les cinq épreuves sont identiques et indépendantes, avec
« succès » = obtenir $B$, de probabilité $\frac12$ :
\[
  \proba{F} = \combi{2}{5}\left(\frac12\right)^{2}\left(\frac12\right)^{3}
  = 10 \times \frac{1}{32} = \frac{5}{16}.
\]
\textbf{b)} On passe par le contraire :
$\proba{X \geq 1} = 1-\proba{X = 0} = 1-\left(\frac12\right)^{5}
= \frac{31}{32}$.
\end{corrige}

% ===========================================================================
\begin{notepourlaclasse}
Le chapitre est réputé facile et il ne l'est pas. Les calculs sont simples,
mais la lecture de l'énoncé est piégeuse : avec ou sans remise, simultané ou
successif, exactement ou au moins. Consacrez la première séance entière à
cette lecture, sans calculer quoi que ce soit — seulement décrire l'univers et
compter.

L'arbre doit être imposé, même quand il paraît superflu. Un élève qui trace
son arbre ne se trompe pas de formule ; un élève qui calcule de tête confond
$\probacond{B}{A}$ et $\proba{A \cap B}$ une fois sur deux.

Enfin, la question finale du sujet — trouver le plus petit $n$ tel que
$p_{n} \geq 0{,}99$ — est un exercice de logarithme autant que de
probabilités. Rappelez que diviser par $\ln$ d'un nombre inférieur à $1$
change le sens de l'inégalité : c'est là que se perdent les points, année
après année.
\end{notepourlaclasse}
